\chapter{State of the Art and Related Work}

\section{Existing Orchestration Frameworks}
ROS creates a separation of concerns by using nodes that execute clearly defined tasks and communicate with one 
another through message passing by implementing a publisher subscriber architecture. In a ROS system, nodes communicate
extensively with one another, but they can be abstracted as microservices that are responsible for executing specific 
actions. These microservices usually interact in complex ways that are difficult to handle manually; for this reason, 
technologies such as Kubernetes \cite{verma2015borg} simplify the orchestration of such systems by managing their lifecycles, scaling applications,
and handling errors, thereby creating a fault-tolerant and dynamic system. Although Kubernetes is heavily used in the software domain, it lacks orchestration mechanisms tailored to domains such
as robotics, as it cannot handle ROS-specific concepts such as node graphs, topic compatibility, or sensor algorithm relationships.
RobotKube can improve deployment portability by defining software components that can be containerized using Docker \cite{merkel2014docker} and orchestrated
by Kubernetes, and it can assist with managing the lifecycle of ROS nodes. On the other hand, RobotKube cannot provide mission-driven
orchestration, as the tool has no inherent understanding of the system\'s purpose. Orchestration is achieved through interoperability
with Kubernetes, while the services handled by RobotKube are built using the ROS framework. RobotKube defines robot specific 
resources such as ROS components and initializes the RobotKube Operator which is a controller program. The Operator reacts to 
changes to ROS components and depending on the situation, it can deploy containers, set up networking or restarts components
which previously failed.

ROScribe \cite{roscribe} uses natural language processing to describe system functionalities and generates the initial ROS graph, 
source code, and launch files. ROScribe builds an entire ROS workspace using four agents that are invoked in a sequential order.
First, the SpecAgent constructs the ROS graph of the project by interpreting user requirements and creates a corresponding
development plan. Second, the GenAgent generates the workspace by either retrieving code from public repositories or 
generating code for the required nodes. Third, the PackAgent analyzes the overall workflow and creates the ROS launch 
files, along with additional support files such as package.xml, CMakeLists.txt, and README.md. In some cases, the generated
ROS workspace may contain errors; in this situation, the SupportAgent is responsible for resolving these issues and ensuring
that the system passes static analysis checks. ROScribe automates code writing, launch file generation, and static code analysis
as part of an offline workspace synthesis process. Its operation is limited to the design phase, generating an initial ROS 
workspace based on user descriptions, without adapting configurations at runtime. Also, ROScribe does not handle
parameter configuration, forcing users to manually create and populate the YAML files associated with each node. This 
limitation is significant in ROS-based systems, where parameters are often tightly coupled to sensors, algorithms, and 
mission requirements, making manual configuration repetitive and error-prone. As a result, ROScribe cannot infer or adapt
configurations based on high-level mission intent, relying instead on static, user-defined parameter settings.

FogROS \cite{fogros22023} is an adaptive framework designed to automate the deployment of fog robotics applications. As robotic algorithms become
increasingly complex, particularly in areas such as perception and mapping, edge-based systems may struggle to handle the 
associated computational load under certain conditions. One approach to addressing this limitation is to leverage cloud 
resources, which can provide scalable computational capabilities for demanding operations, subject to network and latency
constraints. FogROS builds on top of the ROS framework and allows users to specify which components of a ROS computation
graph should be executed on cloud resources. FogROS does not provide reasoning or decision-making capabilities, and all 
deployment configurations must be specified manually by the user. As a result, the system cannot infer mission intent or
dynamically adapt configurations based on high-level task descriptions. An additional feature of FogROS is its ability to
automate container creation for cloud deployment, significantly reducing manual setup effort. The framework also supports
flexible networking and pre-deployed nodes, abstracting many aspects of deployment across local devices and cloud environments,
while still requiring user-defined configuration choices.

OSRF Nexus \cite{osrfnexus} is a framework designed to support dependency management and version control within 
ROS-based software ecosystems. It provides mechanisms for organizing, distributing, and resolving
ROS packages and their dependencies across different systems, helping ensure consistency and 
reproducibility of software environments. Nexus operates at the package and dependency level, allowing
developers to manage compatible versions of ROS components without directly modifying application
logic. However, OSRF Nexus focuses exclusively on software dependency resolution 
and does not address system-level configuration or orchestration. It does not generate ROS graphs,
launch files, or parameter configurations, nor does it provide mechanisms for adapting configurations
based on task requirements or mission context. As a result, while Nexus simplifies dependency management
and improves build-time reproducibility, it does not contribute to mission-driven orchestration or 
dynamic configuration synthesis, relying instead on static, user-defined system setups.

Existing tools for ROS-based system development address specific aspects of configuration, 
deployment, or dependency management, but remain limited when considered from a mission-driven
orchestration perspective. ROScribe automates the generation of ROS workspaces, code, and launch 
files at design time, yet produces static artifacts and requires manual parameter configuration, 
lacking support for runtime adaptation or mission-level reasoning. RobotKube improves deployment 
portability and lifecycle management through Kubernetes integration, but relies on pre-defined 
configurations and operates solely at the infrastructure level, without semantic understanding 
of robotic tasks. FogROS focuses on deployment-time placement of ROS nodes across edge and cloud 
resources, assuming pre-existing applications and static configurations, and does not support 
graph synthesis or dynamic reconfiguration. OSRF Nexus enhances dependency and version management,
but does not address configuration generation, orchestration, or task-aware system behavior.

Overall, these approaches improve isolated stages of the ROS workflow while relying on static, 
user-defined configurations. None provide a reasoning layer capable of translating high-level mission
intent into adaptive ROS configurations, highlighting a gap that motivates the proposed research approach.

\begin{table}[H]
\centering
\caption{Comparison of Existing ROS-Oriented Tools}
\label{tab:ros_tools_comparison}
\renewcommand{\arraystretch}{1.3}
\begin{tabular}{|p{3cm}|p{5.5cm}|p{5.5cm}|}
\hline
\textbf{Tool} & \textbf{Provided Capabilities} & \textbf{Limitations} \\
\hline

ROScribe &
Natural-language-driven generation of ROS workspaces, including source code, ROS graphs, launch files, and basic static analysis using a multi-agent pipeline. &
Operates only at design time; generates static artifacts; does not handle parameter configuration, runtime adaptation, or mission-level reasoning. \\
\hline

RobotKube &
Containerized deployment and lifecycle management of ROS-based components using Kubernetes; improves portability, scalability, and infrastructure-level fault tolerance. &
Relies on pre-defined configurations; lacks semantic understanding of ROS graphs, sensor--algorithm relationships, and mission-driven orchestration. \\
\hline

FogROS &
Deployment-time placement of ROS nodes across edge and cloud resources; automated container creation and infrastructure abstraction for fog robotics applications. &
Assumes pre-existing ROS applications and static configurations; does not support graph synthesis, parameter abstraction, or runtime mission adaptation. \\
\hline

OSRF Nexus &
Dependency and version management for ROS packages; improves reproducibility and consistency of ROS software environments. &
Operates solely at the package and dependency level; does not address configuration generation, orchestration, or task-aware system behavior. \\
\hline

\end{tabular}
\end{table}

\section{Language-Based and AI-Driven Robotics Interfaces}
Despite the fact that the outputs of large language models (LLMs) remain nondeterministic, they can still
be integrated into robotic systems in a controlled manner. One example of such integration is ROScribe, 
which leverages LLM-based workflows to automate the generation of ROS workspaces. Similar to the approach 
described in the AutoGen \cite{wu2023autogen} framework, ROScribe employs a multi-agent pipeline in which agents exchange intermediate
artifacts—such as generated code—and invoke external tools to enrich, validate, or evaluate the produced content.
Another relevant work for the current context \cite{ahn2022saycan} demonstrates how high-level task intent
can be extracted from a single natural language instruction and decomposed into multiple subtasks. In this approach, the 
language model proposes candidate actions expressed in natural language, while the robot provides perception and actuation 
capabilities, effectively serving as the interface to the physical environment. The robot is equipped with a set of predefined 
elementary skills, typically involving both perception and motion, which can be composed to execute more complex tasks.
In this framework, the LLM is used to evaluate the suitability of available skills at each step by estimating their contribution 
to task completion through an affordance-based scoring mechanism. The task is defined using a natural language instruction, 
while each elementary skill is associated with a short description and execution constraints. A similar abstraction could be 
applied to ROS-based systems, where nodes are treated as modular components characterized by their code, parameters, dependencies,
and functional descriptions. By maintaining a database of such components, an LLM could identify suitable building blocks and decompose
a high-level task into lower-level configuration requirements, enabling the mapping of individual components to each subtask.
While these approaches focus on task execution and skill selection, they do not address system-level configuration synthesis, 
which remains the focus of the proposed work.

Large language models have consistently demonstrated the ability to generate and adhere to structured output formats such as JSON \cite{thinkinsidejson2025},
which is a crucial requirement for integrating LLMs into software systems. When developing applications that rely on automated 
reasoning, it is often necessary to convert unstructured inputs, such as natural language descriptions, into a fixed set of 
structured representations that can be programmatically processed. In this context, LLMs do not act as code executors, but rather 
as planners and selectors that identify the appropriate tools, components, or configuration options required to fulfill a given task.
This capability is particularly relevant for configuration synthesis tasks,
where downstream systems depend on well-defined and machine-readable representations, such as parameter sets or orchestration descriptors.

Another foundational contribution to reliable language-based system integration is presented in this paper \cite{ouyang2022instructgpt}. This work demonstrates that large language models can be explicitly trained to prioritize instruction adherence over 
unconstrained text generation through supervised fine-tuning and reinforcement learning from human feedback. By optimizing models to 
follow user-provided instructions and constraints, the approach significantly improves controllability and predictability of generated 
outputs. This capability is essential for applications where natural language serves as a formal interface rather than a conversational medium.
In the context of configuration synthesis, instruction-tuned models are better suited to interpret high-level mission descriptions and translate
them into constrained, structured representations without deviating from specified requirements. When combined with schema-constrained 
generation strategies, instruction-following behavior enables the use of LLMs as reliable planners that map intent to machine-readable 
configuration contexts.

Before modifying model parameters through fine-tuning or reinforcement learning, large language models can be adapted to specific tasks
through prompt-based techniques. These approaches aim to improve model behavior by optimizing how tasks, constraints, and output formats 
are expressed to the model, while keeping the underlying model weights unchanged. Prompt-based adaptation is particularly attractive in 
early-stage system development due to its low computational cost, reversibility, and rapid iteration cycle.
One class of approaches relies on text-based prompt optimization, where prompts remain human-readable and are refined through structured 
experimentation. Libraries such as DSPy \cite{khattab2023dspy} support this paradigm by treating prompts as composable and optimizable programs. DSPy enables the 
construction of modular prompting pipelines and supports evaluation-driven prompt refinement, allowing developers to improve reliability and 
schema adherence without retraining the model. This makes text-based prompt optimization suitable for exploring task formulations and enforcing 
structured outputs in configuration synthesis workflows.
A complementary approach involves soft or continuous prompt tuning, where prompts are represented as learned embedding vectors rather than 
explicit text. Frameworks such as PEFT provide implementations of prompt tuning and prefix tuning techniques, in which a small number of trainable
prompt embeddings are optimized while the base model remains frozen. These methods allow task-specific behavior to be encoded numerically, offering
greater stability and expressiveness than manual prompting, while remaining significantly more efficient than full model fine-tuning.
Together, these approaches provide a progressive path for adapting large language models, enabling reliable task execution through prompt design
before committing to more costly model modification techniques.

Recent advances in large language models have demonstrated strong capabilities in code-aware reasoning and structured program generation. Models 
such as DeepSeek-Coder are explicitly trained on large-scale code corpora and optimized to generate syntactically correct and semantically coherent
source code across multiple programming languages. These models have shown improved performance in tasks such as code completion, refactoring, and 
generation of structured artifacts, making them suitable candidates for applications that require reliable synthesis of machine-readable outputs.
The relevance of code-oriented language models extends beyond traditional software development to configuration-centric domains, where structured 
representations such as JSON or YAML must be generated with high precision. In this context, language models can be leveraged to transform natural
language descriptions into configuration templates, parameter sets, or orchestration descriptors that can be further processed by downstream systems.
This aligns with recent work demonstrating that code-specialized models exhibit stronger adherence to formal structure compared to general-purpose 
language models. However, despite their strengths in program synthesis, models such as DeepSeek-Coder operate without domain-specific knowledge of 
robotic systems. They lack an inherent understanding of ROS concepts, mission semantics, or the relationships between sensors, algorithms, and 
execution contexts. As a result, while code-aware language models can reliably generate structured configuration artifacts, they require an external
abstraction layer to provide semantic grounding and ensure that generated outputs correspond to valid and meaningful robotic system configurations.